\documentclass[10pt,a4paper,sans]{moderncv}
\moderncvstyle{classic}
\moderncvcolor{blue}
\usepackage[utf8]{inputenc}
\usepackage[scale=0.85]{geometry}
\usepackage{helvet}

\name{Alexandro}{Disla}
\title{Gestion de l'Information \& MEAL}
\address{Rue Saint-Louis Jeanty \#14, Route de Frère}{Port-au-Prince, Haïti}{}
\phone[mobile]{(+509) 4148-3700}
\email{alexandrodisla@hotmail.com}
\extrainfo{GitHub: \url{https://github.com/AD0791}}

\begin{document}
\makecvtitle

\section{Profil Professionnel}
\cvitem{}{Spécialiste de la gestion de l'information et du MEAL avec \textbf{cinq ans d'expérience en ONG} en Haïti, sur trois secteurs --- santé et VIH à la Caris Foundation International, eau, hygiène et assainissement chez HANWASH, éducation chez Anseye Pou Ayiti --- précédés de huit ans de suivi d'indicateurs de projets publics au Ministère de la Planification. Je conçois la chaîne entière d'un système d'information de programme : formulaires \textbf{XLSForm} sur \textbf{KoboToolbox} et \textbf{ODK}, base de données des participants qui reçoit les soumissions, contrôles de qualité et dédoublonnage, puis analyse et tableaux de bord pour la décision et le rapportage bailleur. \textbf{Ingénieur de données} en parallèle, je construis et administre ces bases moi-même au lieu de seulement les spécifier. Trois ans sur les données d'un programme VIH m'ont appris à traiter l'information sensible avec la discrétion qu'elle exige. \textbf{Créole haïtien et français maternels}, anglais courant. Résidant à Route de Frère, à proximité de Pétion-Ville.}

\section{Compétences Clés}
\cvitem{}{
\begin{itemize}
    \item \textbf{Systèmes d'information et bases de données :} Conception et administration de bases relationnelles des participants et des activités (MySQL, PostgreSQL) alimentées automatiquement par la collecte ; tableaux de suivi des indicateurs (ITT) ; clé d'identification unique des bénéficiaires, enregistrement, dédoublonnage, mise à jour et sauvegarde des listes.
    \item \textbf{Collecte numérique :} Formulaires XLSForm sur KoboToolbox, ODK, CommCare et mWater --- logiques de saut, contraintes de validation, listes en cascade, codification, collecte hors ligne ; administration du serveur et du panneau d'administration ODK ; réception des soumissions par webhook et intégration directe en base.
    \item \textbf{Qualité des données :} Protocoles \textbf{DQA} établis et supervisés ; contrôles automatisés de complétude et de cohérence, audits d'enregistrements orphelins, nettoyage ; vérification des chiffres rapportés jusqu'aux registres sources et suivi des actions correctives jusqu'à clôture.
    \item \textbf{Protection et sécurité des données :} Données de bénéficiaires d'un programme VIH ; contrôle d'accès par rôle, sessions sécurisées, stockage sécurisé et sauvegardes complètes de la base ; guides techniques documentant les systèmes et leurs procédures.
    \item \textbf{Analyse et visualisation :} \textbf{Excel (avancé)}, \textbf{Power BI}, \textbf{R}, SQL et Python ; statistique descriptive et inférentielle, plans d'échantillonnage (économiste appliqué de formation) ; tableaux de bord de suivi des indicateurs et rapports bailleur vérifiés contre leurs sources.
    \item \textbf{Formation et encadrement technique :} Formation d'agents de collecte et d'équipes partenaires aux outils numériques et aux procédures de qualité ; direction technique des agents de collecte d'un programme multi-sites.
\end{itemize}}

\section{Expériences en Gestion de l'Information et MEAL}

\cventry{Jan 2026 -- Mars 2026}{Analyste de Données \& Consultant MEAL}{Anseye Pou Ayiti (APA)}{Haïti}{}{
\begin{itemize}
    \item \textbf{Système d'information de suivi des bénéficiaires :} Développement et maintenance du back-office de suivi-évaluation de l'organisation --- application web PHP/MySQL en production, alimentée par trois formulaires \textbf{XLSForm} sur ODK/ONA.io dont les soumissions sont reçues par webhook et intégrées automatiquement en base (31 tables).
    \item \textbf{Listes de bénéficiaires :} Clé d'identification unique et déterministe (cohorte, initiales, sexe, année de naissance, séquence) ; pages de \textbf{recherche de doublons}, d'audit des enregistrements orphelins et de réconciliation entre la base de recensement et les fiches de suivi périodique.
    \item \textbf{Intégrité et continuité :} Correction et suppression contrôlées, export Excel, \textbf{sauvegarde complète de la base}, journalisation des soumissions, sessions sécurisées avec verrouillage automatique par inactivité.
    \item \textbf{Indicateurs et formation :} Stratégie de collecte et indicateurs de performance de trois programmes ; administration du système ODK et \textbf{formation des agents de collecte} aux outils et aux procédures de contrôle qualité.
\end{itemize}}

\cventry{Oct 2023 -- Jan 2026}{Consultant MEAL \& Analyste de Données}{HANWASH}{Haïti / Télétravail}{}{
\begin{itemize}
    \item Plans MEAL et tableaux de suivi des indicateurs alignés sur les normes JMP (OMS/UNICEF) pour des programmes d'eau, d'hygiène et d'assainissement.
    \item Enquêtes géospatiales sur mWater et administration de la base ; tableaux de bord automatisés (mWater, \textbf{Power BI}) connectés aux API de collecte ; guides techniques et formation des équipes à leur usage autonome.
\end{itemize}}

\cventry{Oct 2021 -- Août 2024}{Officier de Suivi \& Évaluation}{Impact Youth Project, Caris Foundation International}{Haïti}{}{
\begin{itemize}
    \item \textbf{Système multi-sites :} Cycle de suivi d'un projet de santé et de VIH auprès de jeunes sur des cliniques multi-sites ; standardisation des formulaires et des flux de données entre sites, consolidation des remontées dans les délais.
    \item \textbf{Systèmes d'information :} Mise en place et administration des systèmes de collecte numérique et de la base alimentant le rapportage (CommCare, HIVHaiti, MySQL), sous les exigences de confidentialité propres aux données VIH.
    \item \textbf{Encadrement technique :} Direction technique des agents de collecte sur l'ensemble des sites --- répartition du travail, contrôle qualité, formation aux procédures --- et référent du personnel de suivi-évaluation sur les contrôles de cohérence.
    \item \textbf{Qualité des données (DQA) :} Protocoles de Data Quality Assessment établis et supervisés, chiffres rapportés remontés jusqu'aux registres sources, actions correctives suivies jusqu'à clôture.
    \item \textbf{Rapportage bailleur :} Production et soumission des indicateurs MER destinés à l'USAID via DATIM, la plateforme du PEPFAR bâtie sur DHIS2.
    \item \textbf{Analyse \& restitution :} Rapports d'activité automatisés en Python et SQL, réduisant de 40 \% le temps de traitement manuel ; tableaux de bord \textbf{Power BI} et Quarto pour le suivi périodique des indicateurs.
\end{itemize}}

\cventry{Nov 2015 -- Août 2023}{Économiste --- Direction de Planification Économique et Sociale (DPES)}{Ministère de la Planification et de la Coopération Externe (MPCE)}{Port-au-Prince}{}{Suivi de l'exécution du Plan Triennal d'Investissement 2014--2016 : indicateurs de projets publics sur l'ensemble des secteurs et des départements géographiques ; unité statistique et modélisation, avec participation à la construction du << Modèle de Planification du Développement >>.}

\section{Expériences en Ingénierie de Données et Systèmes}

\cventry{Févr 2026 -- Présent}{Ingénieur Logiciel (Fullstack)}{Tekkod LLC}{Télétravail}{}{Modernisation d'infrastructure mobile avec fiabilité hors ligne et sécurité des données ; tableaux de bord Looker Studio.}

\cventry{Mars 2021 -- Mai 2024}{Développeur Backend \& Ingénieur de Données}{Tekkod LLC}{Télétravail}{}{Administration de bases relationnelles en accès concurrent ; pipelines ETL (Apache Beam) migrant les données MongoDB vers BigQuery ; services REST sécurisés (FastAPI, JWT) avec contrôle d'accès par rôle.}

\cventry{Nov 2015 -- Mars 2021}{Consultant Logiciel \& Spécialiste Intégration de Données}{CassionSoft (Projet UNOPS)}{Haïti}{}{Intégration de données sur un projet sous contrat des Nations Unies, avec exposition aux procédures et standards documentaires onusiens.}

\section{Formation}
\cvitem{2010--2014}{\textbf{Études Supérieures en Économie Appliquée} --- C.T.P.E.A (Centre de Techniques de Planification et d'Économie Appliquée), Port-au-Prince : cycle de quatre ans après le baccalauréat en économie et statistique appliquées. \textit{Scolarité complétée ; Diplôme d'Études Supérieures en attente de la soutenance du mémoire de sortie --- attestation officielle jointe au dossier.}}
\cvitem{2009--2010}{Baccalauréat (Mention) --- Institution Saint-Louis de Gonzague, Port-au-Prince}
\cvitem{Autoformation}{FlatIron School (Software Engineering), CodeWithMosh, Udemy}

\section{Compétences Techniques}
\cvitem{Collecte}{\textbf{KoboToolbox}, \textbf{ODK (Open Data Kit)}, \textbf{XLSForm}, ONA.io, CommCare, mWater, HIVHaiti, DHIS2/DATIM \textit{(soumission de données)}, protocoles DQA, ITT}
\cvitem{Bureautique \& viz}{\textbf{MS Excel (avancé)}, \textbf{Power BI}, Looker Studio, Word, PowerPoint, GSuite, Teams et Zoom}
\cvitem{Analyse statistique}{\textbf{R} (RMarkdown/Quarto), SQL, Python (Pandas), échantillonnage, analyse quantitative et qualitative}
\cvitem{Bases \& systèmes}{MySQL, PostgreSQL, SQLite, MongoDB, PHP, Docker, Git, AWS, Azure, Google Cloud (BigQuery)}
\cvitem{Langues}{\textbf{Créole haïtien (maternel)}, \textbf{Français (maternel)}, \textbf{Anglais (courant, rédaction de rapports)}}

\end{document}
